\chapter{Gender Affirming Surgery}

\section*{Introduction \& Clinical Indications}
Gender affirming surgery encompasses a diverse array of surgical interventions designed to align an individual's physical characteristics with their deeply felt gender identity. These procedures are critical components of comprehensive gender-affirming care, addressing the profound distress known as gender dysphoria that can arise from an incongruence between one's assigned sex at birth and their internal sense of self. Surgeries range from chest reconstruction (e.g., mastectomy for masculinization or augmentation mammoplasty for feminization) and genital reconstruction (e.g., vaginoplasty, phalloplasty, metoidioplasty) to facial feminization or masculinization, and voice modification procedures. The ultimate goal is to alleviate psychological distress, enhance body congruence, and improve overall quality of life and social integration for transgender and gender-diverse individuals.

\begin{itemize}
  \item Gender Confirmation Surgery (GCS)
  \item Sex Reassignment Surgery (SRS)
  \item Gender-Affirming Surgery (GAS)
  \item Gender Reassignment Surgery
  \item Top Surgery (for chest procedures)
  \item Bottom Surgery (for genital procedures)
\end{itemize}

The fundamental principle of gender affirming surgery is to surgically modify primary and/or secondary sexual characteristics to better reflect an individual's affirmed gender, thereby mitigating gender dysphoria. This involves a complex interplay of reconstructive surgical techniques aimed at creating anatomical forms and functions consistent with the patient's identity. For instance, chest masculinization procedures remove glandular and adipose breast tissue and reshape the chest wall to create a masculine contour, often involving nipple-areola complex repositioning and resizing. Conversely, breast augmentation for feminization uses implants to create a feminine chest profile.

Genital surgeries are highly specialized, utilizing various tissue flaps and grafts to construct neovaginas, neophalluses, or modify existing structures. Vaginoplasty, for example, typically involves creating a vaginal canal using penile skin inversion or intestinal segments, along with clitoroplasty and labiaplasty to form external female genitalia. Phalloplasty involves constructing a neophallus using radial forearm free flap or anterolateral thigh flap, often with subsequent urethral lengthening, scrotoplasty, and testicular implants. The technical goals are not only aesthetic but also functional, aiming for sensation, urinary function, and, where desired, sexual penetrative capacity. The rationale is rooted in the profound psychological and social benefits derived from achieving physical congruence, which significantly improves mental health, body image, and overall quality of life for transgender and gender-diverse individuals.

The origins of modern gender affirming surgery can be traced back to the early 20th century, with pioneering efforts in Europe. One of the earliest documented cases of male-to-female genital surgery was performed by Dr. Magnus Hirschfeld and Dr. Ludwig Levy-Lenz in Berlin in 1931 on Dora Richter. Later, Dr. Georges Burou in Casablanca, Morocco, refined vaginoplasty techniques in the 1950s and 1960s, establishing many of the principles still used today, particularly the penile inversion technique. In the United States, Dr. Harry Benjamin's work on gender identity and the establishment of the Harry Benjamin International Gender Dysphoria Association (now the World Professional Association for Transgender Health, WPATH) in 1979 were pivotal in standardizing care and promoting a multidisciplinary approach.

The evolution of techniques has been marked by significant advancements. Early procedures were often rudimentary, focusing primarily on external appearance. Over time, surgical innovations, particularly in microsurgery and tissue transfer, have allowed for more sophisticated and functional reconstructions, especially in phalloplasty and vaginoplasty, improving sensory and aesthetic outcomes. The development of free flap techniques in the 1970s and 80s revolutionized phalloplasty, enabling the creation of larger, sensate neophalluses. The shift from a purely medical model to one that emphasizes patient autonomy, informed consent, and a holistic understanding of gender identity has also profoundly shaped the field, moving from "sex reassignment" to "gender affirmation" and focusing on patient-centered care.

Gender affirming surgery is indicated for individuals who meet specific criteria, primarily driven by the presence of persistent and well-documented gender dysphoria, as outlined by the World Professional Association for Transgender Health (WPATH) Standards of Care.

\begin{itemize}
  \item To alleviate severe and persistent gender dysphoria that significantly impacts mental health and quality of life.
  \item To achieve physical congruence with one's gender identity, reducing the mismatch between internal sense of self and external bodily characteristics.
  \item When non-surgical interventions, such as gender-affirming hormone therapy and psychotherapy, have been insufficient to fully resolve gender dysphoria.
  \item To improve social functioning and reduce anxiety or depression related to gender incongruence.
  \item To enhance self-esteem and body image, leading to improved psychological well-being.
  \item To facilitate passing in one's affirmed gender, which can reduce experiences of discrimination and violence.
  \item As part of a comprehensive, individualized treatment plan developed in consultation with mental health professionals and medical providers.
  \item To address specific physical characteristics that cause distress, such as breast tissue in transmasculine individuals or lack of breast development in transfeminine individuals.
\end{itemize}

Gender affirming surgery is considered a primary treatment modality for gender dysphoria, meaning it directly addresses the core issue of physical incongruence with gender identity. However, it is rarely the initial intervention in a patient's gender affirmation journey. Typically, it is undertaken after a thorough period of psychological assessment, often including psychotherapy, and usually after a sustained period of gender-affirming hormone therapy (GAHT). GAHT can induce significant physical changes (e.g., breast development, voice deepening, fat redistribution) that may reduce the need for certain surgeries or optimize the results of others. For example, chest masculinization may be more effective after a period of testosterone therapy, and breast augmentation is often performed after estrogen therapy has maximized natural breast development. Therefore, while primary in its goal of physical affirmation, surgery functions as a later-stage, definitive intervention within a comprehensive, multidisciplinary care pathway, building upon foundational medical and mental health support.

\section*{Relevant Surgical Anatomy}
Gender affirming surgeries involve a wide range of anatomical regions, each with specific considerations for successful outcomes and complication avoidance. For chest masculinization, the primary structures are the mammary glands, overlying skin and subcutaneous fat, and the underlying pectoralis major muscle. Key neurovascular structures include the lateral thoracic artery and vein, internal mammary artery and vein, and the intercostal nerves, which provide sensation to the chest wall and nipple-areola complex. Preservation of these nerves is crucial for maintaining sensation where possible. Surgical landmarks include the inframammary fold and the sternal notch, guiding the creation of a masculine chest contour. For breast augmentation, the pectoralis major muscle and its fascial attachments are important for submuscular implant placement, and the inframammary fold is a key landmark for incision planning.

Genital surgeries involve complex perineal and pelvic anatomy. Vaginoplasty requires careful dissection within the rectovesical space to create a neovaginal canal, necessitating an intimate understanding of the rectum, bladder, and urethra to prevent injury. The pudendal nerve and its branches are critical for clitoral and labial sensation. Arterial supply to the neovagina often relies on local perineal vessels or, in cases of intestinal vaginoplasty, the mesenteric vasculature. Phalloplasty, particularly using free tissue transfer, involves microvascular anastomosis of recipient vessels (e.g., femoral or deep inferior epigastric vessels) to the donor flap's vascular pedicle (e.g., radial artery and venae comitantes for radial forearm free flap, or lateral circumflex femoral artery for anterolateral thigh flap). Sensory innervation of the neophallus is achieved by coapting donor nerves (e.g., median or lateral femoral cutaneous nerve) to recipient nerves (e.g., dorsal clitoral nerve or ilioinguinal nerve). Urethral reconstruction requires meticulous technique to prevent strictures or fistulas, navigating the existing urethra and creating a neourethra from various tissue sources. Facial feminization surgery targets the bony and soft tissue structures of the skull, including the frontal bone, orbital rims, nasal bones, maxilla, mandible, and chin. The facial nerve and its branches, as well as sensory nerves like the supraorbital, infraorbital, and mental nerves, are paramount to identify and protect to avoid motor weakness or sensory deficits.

\section*{Preoperative Workup \& Preparation}
Extensive preoperative preparation is essential for gender affirming surgery to ensure patient safety, optimize outcomes, and meet established guidelines, such as those from WPATH. This comprehensive process involves medical, psychological, and logistical components.

\begin{itemize}
  \item \textbf{Psychological Assessment:} Comprehensive evaluation by qualified mental health professionals, often requiring letters of readiness for surgery, confirming persistent gender dysphoria, capacity for informed consent, and management of co-occurring mental health conditions.
  \item \textbf{Medical Evaluation:} A thorough physical examination, review of medical history, and baseline laboratory tests (e.g., complete blood count, metabolic panel, coagulation studies, hormone levels, HIV/hepatitis screening).
  \item \textbf{Hormone Therapy:} For many procedures, a period of gender-affirming hormone therapy (typically 6-12 months) is recommended or required to induce desired physical changes and optimize surgical results; hormone adjustments may be needed preoperatively.
  \item \textbf{Anesthesia Clearance:} Evaluation by an anesthesiologist to assess fitness for general anesthesia, especially for patients with pre-existing cardiac, pulmonary, or other systemic conditions.
  \item \textbf{Smoking Cessation:} Patients are strongly advised to cease smoking and nicotine use for several weeks to months prior to surgery to minimize wound healing complications, particularly flap necrosis.
  \item \textbf{Medication Adjustment:} Review and adjustment of medications, particularly blood thinners, herbal supplements, and certain psychiatric medications, in consultation with the surgical team.
  \item \textbf{NPO Status:} Strict adherence to "nil per os" (NPO) guidelines, typically no food or drink for 6-8 hours before surgery, to prevent aspiration during anesthesia.
  \item \textbf{Prophylactic Antibiotics:} Administration of prophylactic antibiotics shortly before incision to reduce the risk of surgical site infection.
  \item \textbf{Informed Consent:} Detailed discussion of the procedure, expected outcomes, potential risks, and alternatives, with signed informed consent, ensuring the patient understands the irreversible nature of many procedures.
\end{itemize}

\textbf{Absolute Contraindications:}
\begin{itemize}
  \item Uncontrolled, severe psychiatric illness (e.g., active psychosis, severe untreated depression) that prevents informed decision-making or poses a significant risk to postoperative recovery.
  \item Active substance abuse that compromises patient safety or adherence to postoperative care instructions.
  \item Unrealistic expectations regarding surgical outcomes that cannot be adequately addressed through counseling.
  \item Significant, unmanaged medical comorbidities (e.g., uncontrolled diabetes, severe cardiovascular disease, active malignancy) that pose an unacceptably high surgical or anesthetic risk.
\end{itemize}

\textbf{Relative Contraindications and Precautions:}
\begin{itemize}
  \item Active smoking or nicotine use: Requires cessation for a specified period (e.g., 4-6 weeks) preoperatively due to increased risk of wound complications, flap necrosis, and infection.
  \item Obesity (BMI >30-35 kg/m\textasciicircum2): May increase surgical risks (e.g., DVT, infection, wound dehiscence) and potentially compromise aesthetic outcomes; weight optimization may be recommended.
  \item Certain autoimmune conditions or immunosuppression: May affect wound healing and infection risk, requiring careful medical management and discussion with specialists.
  \item Lack of consistent social support: Can impact adherence to postoperative care and overall recovery, requiring proactive planning.
  \item Inadequate understanding of the irreversible nature of some procedures or the need for lifelong follow-up (e.g., hormone therapy, dilation for vaginoplasty).
  \item Age: While not an absolute contraindication, surgery in adolescents requires strict adherence to WPATH guidelines, including extensive psychological evaluation and parental consent.
\end{itemize}

\section*{Surgical Technique \& Procedure}
The execution of gender affirming surgery is highly individualized and depends entirely on the specific procedure being performed and the patient's goals. All major gender affirming surgeries are performed under general anesthesia in a hospital setting. Patient positioning is crucial and varies by procedure, ensuring optimal surgical access and patient safety.

For a common procedure like **Bilateral Mastectomy with Nipple-Areola Complex Grafting (Chest Masculinization)**:

\begin{enumerate}
  \item  \textbf{Anesthesia and Positioning:} General anesthesia is administered. The patient is typically positioned supine with arms abducted, allowing access to the chest and axillae.
  2.  \textbf{Incisions:} Common incision patterns include the double incision with free nipple grafts (most common for larger chests), periareolar, or keyhole incisions (for smaller chests with good skin elasticity). The double incision involves two horizontal incisions, one above and one below the pectoralis major muscle, connected laterally.
  3.  \textbf{Tissue Excision:} Through these incisions, the surgeon meticulously removes glandular breast tissue and excess skin and fat, sculpting the chest to achieve a masculine contour. Care is taken to preserve underlying muscle and neurovascular structures.
  4.  \textbf{Nipple-Areola Complex (NAC) Grafting:} For double incision, the original NAC is typically removed, resized, and reshaped as a full-thickness skin graft and then reattached in a higher, more masculine position on the newly contoured chest.
  5.  \textbf{Hemostasis and Drains:} Meticulous hemostasis is achieved to minimize bleeding. Surgical drains (e.g., Jackson-Pratt drains) are often placed to remove any accumulating fluid and prevent hematoma or seroma formation.
  6.  \textbf{Closure:} The skin edges are meticulously approximated and closed in layers using absorbable sutures for deeper layers and non-absorbable sutures or surgical glue for the skin, aiming for optimal scar aesthetics.
  7.  \textbf{Dressing and Compression:} A sterile dressing is applied, often followed by a compression garment or binder to reduce swelling and support the new chest contour during initial healing.
\end{enumerate}

Other procedures, such as vaginoplasty or phalloplasty, involve highly complex microvascular and reconstructive techniques, often requiring multiple stages and specialized teams. Vaginoplasty typically involves creating a neovaginal canal, clitoroplasty, and labiaplasty, often using penile skin inversion or peritoneal flaps. Phalloplasty involves constructing a neophallus from a donor site (e.g., radial forearm or anterolateral thigh), urethral lengthening, scrotoplasty, and placement of testicular implants, often in multiple stages. Facial feminization surgery involves osteotomies and contouring of facial bones (e.g., forehead, jaw, chin) and soft tissue modifications to soften masculine features.

\section*{Complications \& Risks}
As with any major surgical procedure, gender affirming surgeries carry general risks, in addition to specific risks related to the complexity and anatomical sites involved.

\begin{itemize}
  \item \textbf{General Surgical Complications:}
    \begin{itemize}
      \item \textbf{Bleeding and Hematoma:} Accumulation of blood under the skin, potentially requiring drainage.
      \item \textbf{Infection:} Surgical site infection, ranging from superficial cellulitis to deep abscess, potentially requiring antibiotics or further surgery.
      \item \textbf{Anesthesia Complications:} Reactions to anesthetic agents, respiratory or cardiac complications, nausea, vomiting, or nerve damage from positioning.
      \item \textbf{Deep Vein Thrombosis (DVT) and Pulmonary Embolism (PE):} Blood clots forming in the legs that can travel to the lungs, a potentially life-threatening complication.
      \item \textbf{Poor Wound Healing/Dehiscence:} Separation of wound edges, especially in areas of tension or compromised blood supply.
      \item \textbf{Adverse Scarring:} Hypertrophic scars or keloids, which may require revision.
    \end{itemize}
  \item \textbf{Procedure-Specific Risks (Examples):}
    \begin{itemize}
      \item \textbf{For Chest Masculinization (Mastectomy):}
        \begin{itemize}
          \item Nipple-areola complex (NAC) necrosis or partial loss.
          \item Changes in NAC sensation (often reduced or absent).
          \item Contour irregularities or asymmetry of the chest.
          \item Seroma formation (fluid collection) requiring aspiration.
          \item Persistent pain or numbness.
        \end{itemize}
      \item \textbf{For Vaginoplasty:}
        \begin{itemize}
          \item Vaginal stenosis (narrowing) requiring lifelong dilation.
          \item Rectal injury or fistula formation (rectovaginal fistula).
          \item Urethral injury or stricture.
          \item Loss of sensation in neoclitoris or neovagina.
          \item Neovaginal prolapse.
          \item Hair growth within the neovagina (if not adequately depilated preoperatively).
        \end{itemize}
      \item \textbf{For Phalloplasty/Metoidioplasty:}
        \begin{itemize}
          \item Urethral complications: stricture, fistula, necrosis of neourethra, requiring multiple revisions.
          \item Flap complications: partial or complete flap necrosis, requiring debridement or re-operation.
          \item Loss of sensation in neophallus.
          \item Testicular implant extrusion or infection.
          \item Erectile dysfunction (if erectile device is implanted).
          \item Donor site morbidity (e.g., scarring, nerve injury, functional deficit).
        \end{itemize}
      \item \textbf{For Facial Feminization Surgery (FFS):}
        \begin{itemize}
          \item Nerve injury (e.g., facial nerve, supraorbital nerve) leading to weakness or numbness.
          \item Hairline recession or alopecia.
          \item Asymmetry.
          \item Bone non-union or infection.
          \item Changes in voice pitch (if laryngeal procedures are performed).
        \end{itemize}
    \end{itemize}
\end{itemize}

\section*{Postoperative Care \& Recovery}
The immediate postoperative period involves close monitoring in the Post-Anesthesia Care Unit (PACU) as the patient recovers from anesthesia. Pain management, nausea control, and monitoring for immediate complications like bleeding are paramount. For chest surgery, patients typically stay overnight or are discharged the same day, while more complex genital surgeries often require a hospital stay of several days to a week. During this initial phase, drains are managed, and patients are encouraged to ambulate early to prevent deep vein thrombosis (DVT) and promote circulation. Initial pain is managed with prescribed oral or intravenous analgesics, and a compression garment is usually applied.

Over the first 1-2 weeks, patients will manage wound care, drain output (if applicable), and pain with prescribed medications. Compression garments are often worn continuously for several weeks to months to reduce swelling and support contouring. Activity restrictions are significant, including avoiding heavy lifting, strenuous exercise, and sometimes specific movements to protect surgical sites. For vaginoplasty, a crucial component of recovery is regular vaginal dilation, which begins shortly after surgery and is a lifelong commitment to maintain depth and prevent stenosis. For phalloplasty, meticulous wound care, especially around the neourethra and donor site, is vital. Long-term recovery involves gradual resumption of normal activities, scar management, and ongoing follow-up with the surgical team and endocrinologist for hormone management. Full healing and maturation of results can take 6-18 months, with sensation potentially returning and improving over an even longer period.

During gender affirming surgery, the surgeon meticulously documents intraoperative findings, which can include the volume and consistency of excised breast tissue, the presence of any anatomical variations, the quality of donor tissue flaps, and the integrity of reconstructed structures. For chest masculinization, findings typically include dense glandular and adipose breast tissue, which is excised to the pectoralis fascia. For vaginoplasty, the depth of the perineal dissection and the integrity of the rectal and urethral walls are critical findings. In phalloplasty, the successful microvascular anastomosis and nerve coaptation are key intraoperative observations.

Pathology reports on resected tissues confirm the nature of the excised material. For mastectomy, the report typically confirms the presence of benign mammary gland tissue, often with fibrocystic changes, and rules out any unexpected malignancy. For genital surgeries, excised tissues (e.g., testes, ovaries, uterus) are sent for pathological examination to confirm their identity and rule out any underlying pathology. Skin grafts or flap tissues are generally not sent for routine pathology unless there is a specific concern. These reports serve as a definitive record of the removed tissues and are crucial for comprehensive patient care and future medical management.

A normal, successful postoperative course following gender affirming surgery is characterized by a predictable progression of healing and resolution of acute symptoms. Initially, patients can expect moderate pain, managed effectively with analgesics, which gradually subsides over days to weeks. Swelling and bruising are common and will progressively decrease over several weeks to months. Wounds should heal without significant infection or dehiscence, with scars maturing and fading over 12-18 months. For chest surgeries, the desired masculine or feminine contour should become evident as swelling resolves. For vaginoplasty, consistent dilation will maintain neovaginal depth and width, and for phalloplasty, successful urethral healing will allow for standing urination. Crucially, a normal course includes a significant reduction in gender dysphoria, improved body image, and enhanced psychological well-being, reflecting the primary goal of the surgery.

Postoperative follow-up is a critical, ongoing aspect of gender affirming care, often extending for many years to ensure optimal long-term outcomes and address any potential complications.

\begin{itemize}
  \item \textbf{Surgical Follow-up Visits:} Regular appointments with the surgical team to monitor wound healing, remove sutures/staples, assess for complications (e.g., seroma, infection), and manage drains.
  \item \textbf{Hormone Management:} Lifelong follow-up with an endocrinologist or primary care provider to monitor hormone levels and adjust therapy as needed, as hormones play a crucial role in maintaining secondary sex characteristics and overall health.
  \item \textbf{Physical Therapy/Rehabilitation:} May be recommended for certain procedures (e.g., chest surgery for range of motion, phalloplasty for mobility or donor site recovery).
  \item \textbf{Dilation Protocol (Vaginoplasty):} Strict adherence to a lifelong dilation regimen is essential to maintain vaginal depth and width and prevent stenosis.
  \item \textbf{Urologic Follow-up (Phalloplasty):} Regular monitoring for urethral complications such as strictures or fistulas, which may require further intervention.
  \item \textbf{Routine Screenings:} Maintenance of appropriate gender-affirming health screenings (e.g., prostate exams for trans women on hormones, breast/cervical cancer screenings for trans men who retain those organs).
  \item \textbf{Revision Procedures:} Discussion and planning for potential secondary or revision surgeries to optimize aesthetic or functional outcomes, which are common for complex reconstructive procedures.
  \item \textbf{Warning Signs:} Patients are educated on warning signs such as fever, increasing pain, redness, swelling, purulent discharge, or sudden changes in sensation, which warrant immediate medical attention.
\end{itemize}
Adherence to these follow-up protocols is paramount for achieving and maintaining the best possible results.

\section*{Outcomes \& Prognosis}
The incidence of gender affirming surgeries has been steadily increasing globally, reflecting greater societal acceptance, improved access to care, and evolving medical guidelines. While precise global statistics are challenging to ascertain due to varying data collection methods, studies from North America and Europe indicate a significant rise in procedures over the past two decades. For instance, data from the United States suggests a substantial increase in the number of gender affirming surgeries performed annually, with chest masculinization and breast augmentation being among the most common procedures. Genital surgeries, while less frequent, also show an upward trend. The age distribution of patients undergoing GAS is broad, ranging from adolescents (with strict protocols) to older adults, though the majority are typically young to middle-aged adults. There is no significant sex predilection for seeking gender affirmation overall, as individuals assigned male and female at birth pursue various procedures to align with their identities. Morbidity rates vary significantly by procedure, with minor complications being relatively common but severe, life-threatening complications being rare. Mortality directly attributable to gender affirming surgery is exceedingly low, comparable to other elective surgeries.

The outcomes of gender affirming surgery are overwhelmingly positive, with high rates of patient satisfaction and significant improvements in mental health and quality of life. Studies consistently report a substantial reduction in gender dysphoria, anxiety, and depression following surgery, with many patients experiencing enhanced self-esteem and social integration. For chest masculinization, satisfaction rates often exceed 90\%, with good aesthetic and functional results. Vaginoplasty typically achieves satisfactory cosmetic and functional outcomes, including sexual sensation and penetrative capacity, though maintaining vaginal depth requires lifelong dilation. Phalloplasty, while technically more complex and associated with higher rates of complications (particularly urethral), also yields high satisfaction regarding body congruence and standing urination, with ongoing advancements continually improving outcomes. Recurrence of gender dysphoria after successful surgery is rare, underscoring the profound and lasting impact of these interventions. Prognostic factors for positive outcomes include comprehensive preoperative psychological assessment, adherence to postoperative care, strong social support, and realistic expectations regarding surgical results and potential need for revisions. While revision surgeries are sometimes necessary to optimize aesthetic or functional outcomes, the overall prognosis for improved well-being and successful gender affirmation is excellent for individuals who undergo these procedures within a multidisciplinary care framework.

\section*{References}
\begin{enumerate}
  \item World Professional Association for Transgender Health (WPATH). Standards of Care for the Health of Transgender and Gender Diverse People, Version 8. WPATH; 2022.
  \item MedlinePlus. Gender Affirming Surgery. U.S. National Library of Medicine; 2024. Available from: \texttt{https://medlineplus.gov/genderaffirmingsurgery.html}
  \item Hage JJ, et al. Gender-Affirming Surgery: Principles and Practice. Thieme; 2020.
  \item American Society of Plastic Surgeons. Gender Confirmation Surgery. Available from: \texttt{https://www.plasticsurgery.org/reconstructive-procedures/gender-confirmation-surgery}
\end{enumerate}

\clearpage
